% title: Coprime-Detecting Function and Prime Support
% short-title: Coprime-Detecting Function
% abstract: We formalize IMO 2024 Shortlist N7. For a positive-integer-valued function satisfying the given coprimality-detecting identity, we prove that for every positive integer \(n\), the possible values of \(f(n)\) are exactly the positive integers with the same set of prime divisors as \(n\).
% subjclass: 03B35
% keywords: olympiad number theory, functional equation, prime support, Lean 4
% author: Mario Hernández M.

\subsection{Coprime-Detecting Function and Prime Support}

\begin{problemstatement}
Let \(\mathbb{Z}_{>0}\) denote the set of positive integers. Let
\(f:\mathbb{Z}_{>0}\to\mathbb{Z}_{>0}\) be a function satisfying the following
property: for \(m,n\in\mathbb{Z}_{>0}\), the equation
\[
  f(mn)^2 = f(m^2)\,f(f(n))\,f\bigl(mf(n)\bigr)
\]
holds if and only if \(m\) and \(n\) are coprime. For each positive integer
\(n\), determine all the possible values of \(f(n)\).
\end{problemstatement}

We formalize the defining predicate
\lean{Biblioteca.Demonstrations.CoprimeDetectingFunction} and prove the
classification theorem
\lean{Biblioteca.Demonstrations.imo2024_sl_n7_values}.

\begin{theorem}
Fix \(n\in \mathbb{Z}_{>0}\). A positive integer \(k\) can occur as \(f(n)\) for
some function \(f\) satisfying the hypothesis if and only if \(k\) and \(n\)
have exactly the same prime divisors.
\end{theorem}

\begin{proof}
We first extract the basic algebraic consequences of the defining identity.
Substituting \(m=n=1\) gives \(f(1)=1\). Then the specializations
\((m,n)=(1,n)\) and \((m,n)=(n,1)\) imply
\[
  f(f(n))=f(n), \qquad f(n^2)=f(n).
\]
In Lean these reductions are
\lean{Biblioteca.Demonstrations.coprimeDetecting_one_eq_one},
\lean{Biblioteca.Demonstrations.coprimeDetecting_idem}, and
\lean{Biblioteca.Demonstrations.coprimeDetecting_sq}. After simplifying the
original hypothesis with these identities, we obtain the equivalent form
\[
  f(mn)^2 = f(m)\,f(n)\,f\bigl(mf(n)\bigr)
  \iff \gcd(m,n)=1,
\]
formalized as \lean{Biblioteca.Demonstrations.coprimeDetecting_Q}, together
with the derived criterion
\[
  f\bigl(mf(n)\bigr)=f(m)f(n)
  \iff \gcd\bigl(m,f(n)\bigr)=1,
\]
formalized as \lean{Biblioteca.Demonstrations.coprimeDetecting_R}.

The first structural step is that every prime divisor of \(n\) also divides
\(f(n)\). Indeed, if a prime \(p\) divides \(n\), write \(n=p^v t\) with
\(v\ge 1\) and \(p\nmid t\). Then \(p^v\) is not coprime to \(f(p^v)\), so
\(p\mid f(p^v)\). Applying the coprime case of the simplified relation to the
pair \((p^v,t)\) yields \(p\mid f(n)\). This is the lemma
\lean{Biblioteca.Demonstrations.coprimeDetecting_prime_dvd_image}.

The key restriction is that for a prime \(p\), the image \(f(p)\) has no prime
divisors other than \(p\). Suppose, toward a contradiction, that a distinct
prime \(q\) divides \(f(p)\). Using the symmetry of the coprimality-transfer
statement
\lean{Biblioteca.Demonstrations.coprimeDetecting_coprime_image_transfer},
formalized further in
\lean{Biblioteca.Demonstrations.coprimeDetecting_prime_dvd_iff_not_coprime_prime_image},
one proves that whenever an integer \(x\) is not coprime to \(pq\), then
\(p\mid f(x)\). Choose \(X\) minimizing the \(p\)-adic valuation of \(f(X)\)
among all such \(x\). Writing
\[
  f(X)=p^m u \qquad (p\nmid u),
\]
one checks that \(q\mid u\), so \(u\) is again not coprime to \(pq\). Applying
the coprime identity to \((p^m,u)\) forces the right-hand side to be divisible
by \(p^{3m}\), while the left-hand side is \(f(X)^2\), whose \(p\)-adic
valuation is exactly \(2m\), a contradiction. The formal contradiction is
\lean{Biblioteca.Demonstrations.coprimeDetecting_prime_image_unique}.

Now let \(r\) be a prime divisor of \(f(n)\). If \(r\nmid n\), then the
criterion
\lean{Biblioteca.Demonstrations.coprimeDetecting_prime_dvd_iff_not_coprime_prime_image}
shows that \(n\) is not coprime to \(f(r)\). But every prime divisor of
\(f(r)\) must equal \(r\), contrary to \(r\nmid n\). Hence every prime divisor
of \(f(n)\) already divides \(n\). Combined with the previous inclusion, this
gives equality of prime supports:
\[
  \operatorname{PrimeDiv}(f(n))=\operatorname{PrimeDiv}(n).
\]
This is the theorem
\lean{Biblioteca.Demonstrations.coprimeDetecting_primeFactors_image}.

For the converse, fix a positive integer \(k\) with the same prime divisors as
\(n\). Define
\[
  g_k(m)=\prod_{p\mid m} p^{e_p},
  \qquad
  e_p=
  \begin{cases}
    v_p(k), & p\mid k,\\
    1, & p\nmid k.
  \end{cases}
\]
In Lean this family is
\lean{Biblioteca.Demonstrations.primeSupportModel}. Its prime divisors are
exactly those of the input, and the product identity over prime-factor sets
shows that
\[
  g_k(mn)=g_k(m)\,g_k(n)
  \iff \gcd(m,n)=1.
\]
After also observing that \(g_k(g_k(n))=g_k(n)\) and \(g_k(n^2)=g_k(n)\), one
checks directly that \(g_k\) satisfies the original condition; this is
\lean{Biblioteca.Demonstrations.primeSupportModel_isSolution}. Finally, if
\(k\) and \(n\) have the same prime support, then by construction
\[
  g_k(n)=k,
\]
formalized as
\lean{Biblioteca.Demonstrations.primeSupportModel_eq_of_primeFactors_eq}.

Therefore the possible values of \(f(n)\) are precisely the positive integers
having the same set of prime divisors as \(n\).
\end{proof}
